% Related Work — SPIE: section titles ALL CAPS; citations via \cite{}
\section{RELATED WORK}
\label{sec:related}

\rev{In neural rendering, Neural Radiance Fields (NeRF) employ an implicit radiance field representation, allowing high-quality images of unseen viewpoints to be synthesized through differentiable rendering}~\cite{Mildenhall21_NeRF}. Under sparse views, however, insufficient observations cause severe geometric ambiguity. Depth-guided approaches further improve structural consistency~\cite{Wang23_SparseNeRF}, while other methods constrain unobserved regions~\cite{Niemeyer22_RegNeRF}. \rev{In contrast to implicit radiance field representations, 3D Gaussian Splatting (3DGS) models scenes explicitly using optimizable Gaussian primitives, with differentiable rasterization supporting efficient training and real-time rendering}~\cite{Kerbl23_3DGS}. Beyond the original formulation, structured Gaussian representations improve robustness and view-adaptive rendering~\cite{Lu24_ScaffoldGS}. Depth-based approaches introduce additional geometric supervision~\cite{Li24_DNGaussian}, while densification-based methods improve Gaussian coverage~\cite{Zhu24_FSGS}. Other strategies include epipolar depth priors~\cite{Zheng25_NexusGS}. Others explore co-regularization~\cite{Zhang24_CoRGS} and Gaussian dropout~\cite{Park25_DropGaussian}.

\rev{Monocular depth estimators provide useful cues for underconstrained views}~\cite{Yang24_DepthAnythingV2}. Although depth supervision, Gaussian densification, and training-time regularization have improved sparse-view reconstruction, limited-view observations can still result in incomplete initial geometry and overfitting during optimization. To address these challenges, SynGS combines explicit visual-hull constraints with a VGGT-based multi-view depth prior~\cite{Wang25_VGGT} to augment the initial Gaussian geometry before optimization, and incorporates Dynamic Visibility Regularization during training to alleviate sparse-view overfitting.
